\documentclass[dvipsnames,hyperref={pdfpagelabels=false}]{beamer}
% beamer already loads hyperref; do NOT load \usepackage{hyperref} again.

\usepackage{ragged2e} % 文字向兩邊對齊
\renewcommand{\raggedright}{\leftskip=0pt \rightskip=0pt plus 0cm}

\usepackage{listings}
\usepackage{lmodern}
\usepackage{fontspec}
\usepackage{xeCJK}
\setCJKmainfont{Noto Sans CJK TC}
\setCJKsansfont{Noto Sans CJK TC}
\setCJKmonofont{Noto Sans Mono CJK TC}
\usepackage[english]{babel}
\usepackage{amsmath,amssymb,amsfonts,amsthm}

% hyperref settings (since beamer already loads hyperref)
\hypersetup{
  colorlinks=true,
  urlcolor=blue,
  citecolor=blue,
  pdfborder={0 0 0}
}


\usepackage{accsupp}% http://ctan.org/pkg/accsupp
\newcommand{\emptyaccsupp}[1]{\BeginAccSupp{ActualText={}}#1\EndAccSupp{}}


\usepackage{booktabs}
\usepackage{colortbl}
\usepackage{xcolor}
\usepackage{threeparttable}

\usepackage{tikz}
\usetikzlibrary{arrows,decorations.pathmorphing,backgrounds,fit,positioning,shapes.symbols,chains}
\usetikzlibrary{tikzmark}

\newcommand{\indep}{\perp\!\!\!\perp}
\newcommand{\nindep}{\not\!\perp\!\!\!\perp}

\beamertemplatefootpagenumber



\title[]  
{Guidelines for Term Paper}

%\subtitle
%{免除部分負擔對幼兒醫療利用與健康的影響:斷點迴歸分析}

%\author 
 %{Tzu-Ting Yang\inst{1} \and Hsing-Wen Han\inst{2} \and Hsien-Ming Lien\inst{3}}

%\institute[short]{\inst{1} Academia Sinica \and \inst{2} Tamkang University \and \inst{3} National Chengchi University}


\author 
{Prof. Tzu-Ting Yang  \\ 楊子霆}

%\institute{Institute of Economics, Academia Sinica  \\ Econ, NTU}
\institute{Institute of Economics, Academia Sinica  \\ 中央研究院經濟研究所}
%\institute{Institute of Economics, Academia Sinica  \\ Econ, NTU}
%\institute{Institute of Economics, Academia Sinica  \\ IMES/IDAS, NCCU}

%\author 
%{楊子霆 \\ 中研院經濟所}

%\institute 
%{中正大學國際經濟專題研討}




\date
{\today}


%\usetheme{Copenhagen}  %applied
%\usetheme{CambridgeUS}  %labor
%\usepackage{beamerthemeshadow}  % data (commented: use default beamer theme)
%\usecolortheme{seahorse}
\setbeamertemplate{footline}{%
	\hfill\usebeamercolor[fg]{page number in head/foot}%
	\usebeamerfont{page number in head/foot}%
	\insertframenumber\,/\,\inserttotalframenumber\hspace{1em}\vspace{2pt}%
}

%\usepackage{beamerthemeshadow}
\beamersetuncovermixins{\opaqueness<1>{25}}{\opaqueness<2->{15}}

\begin{document}


\begin{frame}{}
 \titlepage
\end{frame}


\begin{frame}{Term Paper Deadline: 6/18}
\begin{itemize}
	
	\item You should submit your term paper before 6/18
	\vspace{0.2cm} 
	\begin{itemize}
		\item No late submission 
		\vspace{0.2cm}
		\item Let me know if you want me to submit grade earlier
	\end{itemize}
	\vspace{0.2cm}
	\item Write your paper in Chinese or English
   \vspace{0.2cm}
    \item You need to type your paper using \textbf{Latex}
	
\end{itemize}
\end{frame}



\begin{frame}{Term Paper Deadline: 6/18}
	\begin{itemize}
		
	  \item \textbf{Please upload your paper before 11:59pm 6/18 using the following link}
	\vspace{0.2cm}
	\begin{itemize}
		\item \url{https://www.dropbox.com/request/aBeYTc1TAe2LUXbRj4VA}
		\vspace{0.2cm}
		\item Format of file name: studentID+name
		\vspace{0.2cm}
		\item Example: 1095323010\_TzuTingYang
		\vspace{0.2cm}
	\end{itemize}
		
	\end{itemize}
\end{frame}



\begin{frame}{Term Paper Deadline: 6/18}
\begin{itemize}
	
\item Letter style: roughly 2,000-3,000 words (less than 10 pages)
\vspace{0.3cm}
\item The maximum number of exhibits (figures + tables) is five
%6-10 pages including tables, figures, footnotes, appendices, and references
\vspace{0.3cm}
\begin{itemize} 
	\item See \textbf{Economics letters}
	\vspace{0.3cm}
	\item \url{https://www.sciencedirect.com/journal/economics-letters}
		\vspace{0.3cm}
	\item \url{https://www.sciencedirect.com/science/article/pii/S0165176524001617}
	
			\vspace{0.3cm}
	
	\item \url{https://www.sciencedirect.com/science/article/pii/S0165176524001836}
	
\end{itemize}
\vspace{0.3cm}
	
\end{itemize}
\end{frame}



\begin{frame}{Guideline for Writing a Term Paper}
\begin{itemize} 
	
	\item Use credible causal inference methods to answer an empirical question
	\vspace{0.3cm}
	\begin{itemize} 
		\item Test economics (social science) theory
		\vspace{0.3cm}
		\item Estimate policy effect
		\vspace{0.3cm}
		\item Any interesting questions regarding to human behavior/social phenomenon
	\end{itemize}
	
	
\end{itemize}
\end{frame}



%\item Co-authorship is allowed, but comes with higher expectations
%\vspace{0.3cm}
%\item At most two authors
%\vspace{0.3cm}
%\item Discuss your research question (3/22) and empirical strategy (4/26) with me
%\vspace{0.3cm}
%\item Presentation (6/21) and final paper (7/31)
%\end{itemize}
%\end{frame}



\begin{frame}{Guideline for Writing a Term Paper}
\begin{itemize} 
\item The Typical Structure of an Empirical Paper
\vspace{0.3cm}
\begin{itemize} 
\item[1] \textbf{Introduction/Motivation}
\vspace{0.3cm}
\item[2] Literature Review
\vspace{0.3cm}
\item[3] Theoretical Framework 
\vspace{0.3cm}
\item[4] \textbf{Data and Sample}
\vspace{0.3cm}
\item[5] \textbf{Empirical Method}
\vspace{0.3cm}
\item[6] \textbf{Results} 
\vspace{0.3cm}
\item[7] \textbf{Discussion and Conclusion}
\end{itemize}

\end{itemize}
\end{frame}




\begin{frame}{Guideline for Writing a Term Paper}
	\begin{itemize} 
		\item[0] Title of the paper
		\vspace{0.3cm}
		\begin{itemize} 
			
			\item You need to give the title of your paper
					\vspace{0.3cm}
			\item Use one sentence to summarize the research question you want to answer in this paper
			
			\vspace{0.3cm}
			\begin{itemize} 
			\item 失業對健康的影響
						\vspace{0.3cm}
			\item The Effect of Financial Resource on Fertility: Evidence from Lottery Winners
					\end{itemize}
		\end{itemize}
		
		
	\end{itemize}
\end{frame}



\begin{frame}{1. Introduction/Motivation}
	\begin{itemize} 
		
		\item The question you are trying to address
		\vspace{0.2cm}
		\begin{itemize} 
			\item Stating the hypothesis to be tested directly is a good way to do this
		\end{itemize}
		\vspace{0.2cm}
		\item Briefly state why we should care about this question 
		\vspace{0.2cm}
		\begin{itemize}
			\item Is it an unproven theoretical result?
			\vspace{0.2cm} 
			\item An important policy question or an interesting human behavior/social phenomenon?
		\end{itemize}
		\vspace{0.2cm}
		\item Briefly state how you answer this question 
		\vspace{0.2cm}
		\begin{itemize}
			\item Empirical strategy you use to answer the question
			\vspace{0.2cm} 
			\item What data you are using
			\vspace{0.2cm}
		\end{itemize}
		\vspace{0.2cm}
		\item Briefly discuss your main findings  and contributions
		\vspace{0.2cm}
		\item \textbf{We usually finish ``Introduction" in the end of project}
		
		
		
	\end{itemize}
	
	
	
\end{frame}




\begin{frame}{2. Literature Review}
	\begin{itemize} 
		
		\item The first section should discuss previous research that is \textbf{directly relevant to your paper}
		\vspace{0.2cm}
		\begin{itemize} 
			\item Not every single paper written on the topic
		\end{itemize}
		\vspace{0.2cm}
		\item The second section should explain your contribution in more detail
		\vspace{0.2cm}
		\item You should discuss how your approach is different from what has been done before
		\vspace{0.2cm}
		\begin{itemize}
			\item Answer a question more broadly/specifically
			\vspace{0.2cm}
			\item Use different data
			\vspace{0.2cm} 
			\item Use different empirical strategy
			\vspace{0.2cm} 
			\item Comparing how you are improving on a previous paper is useful
		\end{itemize}
		\vspace{0.2cm}
		\item You can combine introduction with literature review when you mention your contributions.
		
		
	\end{itemize}
	
	
	
\end{frame}





\begin{frame}{3. Theoretical Framework}
	\begin{itemize} 
		
		\item If your paper does not try to verify economics theory, you might not need this part
		\vspace{0.2cm}
		\item Not every paper will require the development of a full theoretical model
		\vspace{0.2cm}
		\item In the end, you need an empirical model, so the theoretical model you develop must lead somehow to what you are testing
		
		
		
	\end{itemize}
	
	
	
\end{frame}




\begin{frame}{4. Data and Sample }{Data}
	\begin{itemize} 
		
		\item Describe the name and source of the data you are using and the period it covers
		\vspace{0.2cm}
		\item Describe whether you have a panel, cross section or
		time series, what the unit of observation is and how many observations you have
		\vspace{0.2cm}
		\item You may want to highlight the important limitations of this data
		
		
		
	\end{itemize}
	
	
	
\end{frame}



\begin{frame}{4. Data and Sample }{Sample}
	\begin{itemize} 
		
		\item The second section should present (relevant) descriptive statistics of the estimated sample
		\vspace{0.2cm}
		\item You should have a couple of tables with means and standard deviations for the variables you will be using in the analysis
		\vspace{0.2cm}
		\item You may want to present these descriptive statistics
		for different subgroups
		
	\end{itemize}
	
	
	
\end{frame}






\begin{frame}{5. Empirical Method}
	\begin{itemize} 
		
		\item You should write down the basic econometric specification first and explain each of the variables and the parameters of interest
		\vspace{0.2cm}
		\begin{itemize} 
			\item The model could be OLS, PDS, matching, IV, DID, RDD, SCM or others
			\vspace{0.2cm}
			\item Since IV, DID, RDD, and SCM can deal with unobserved confounding factors, these methods usually give us more credible results
		\end{itemize}
		\vspace{0.2cm}
		\item Discuss the identification assumption of your method.
		\vspace{0.2cm}
		\item Discuss why this is the correct/best specification for the
		question you wish to address.
		\vspace{0.2cm}
		\item Discuss why you want to control for certain variables
		
		%certain variables are included and others not.
		
	\end{itemize}
	
	
	
\end{frame}





\begin{frame}{6. Results}
	\begin{itemize} 
		
		\item You should present results in a way that develops your argument step-by-step
		\vspace{0.1cm}
		\begin{itemize} 
			\item For example, 1. main results; 2. break those results down by subgroups; 3. perform robustness checks
		\end{itemize}
		\vspace{0.1cm} 
		\item Any tables with parameter estimates should clearly state which outcome variable you are using
		\vspace{0.1cm}
		\item You should let readers know which control variables are included in your empirical models (put this information in table notes)
		\vspace{0.1cm}
		\item Just discuss the most interesting and important
		estimates in your discussion of the table: \textbf{estimated causal effects}
		\vspace{0.1cm}
		\item Make sure you report standard errors with your estimates in the tables
		\vspace{0.1cm}
		\item Just look at some economics journals for a good table
		format
		
	\end{itemize}
	
\end{frame}



\begin{frame}{6. Results}
	\begin{itemize} 
		
		\item Interpret the magnitude of your parameter estimates in an economically meaningful way
		\vspace{0.2cm}
		\begin{itemize} 
			\item For example, ``we find that b=0.3, so that increasing X by
			one unit increases y by 0.3. The implied elasticity is..."
		\end{itemize}
		\vspace{0.2cm} 
		
		\item \textbf{Graphical evidence is a very convincing way to show your reader that you have found something real}
		\begin{itemize} 
			\vspace{0.2cm}
			\item Graphs are worth a thousand words
		\end{itemize}
		\vspace{0.2cm}
		\item Discuss whether the key estimates are statistically significant
		\vspace{0.2cm}
		\begin{itemize} 
			\item Compare your results to what others have found
			\vspace{0.2cm}
		\end{itemize}
		\item \textbf{Don't worry if you don't find anything significant as long as your methods are credible and you have interpreted the results well}
		
	\end{itemize}
	
\end{frame}







\begin{frame}{7. Discussion and Conclusion }
	\begin{itemize} 
		
		\item Briefly summarize your findings
		\vspace{0.2cm}
		\item Discuss some limitations in this paper 
		\vspace{0.2cm}
		\item Discuss your future plan about this project if you have more time and data
	\end{itemize}
	
\end{frame}



%
%\title[]  
%{Utilize AI as Your Research Assistant}
%
%
%\author 
%{}
%
%%\institute{Institute of Economics, Academia Sinica  \\ Econ, NTU}
%%\institute{Institute of Economics, Academia Sinica  \\ IMES/IDAS, NCCU}
%\institute{}
%
%%\author 
%%{楊子霆 \\ 中研院經濟所}
%
%%\institute 
%%{中正大學國際經濟專題研討}
%
%
%
%
%\date
%{}
%
%
%\begin{frame}{}
%	\titlepage
%\end{frame}
%
%
%
%
%\begin{frame}{Utilize AI as Your Research Assistant}
%	\begin{itemize} 
%		\item You should use the available AI tools to help you conduct research
%		\vspace{0.2cm}
%		\item Utilizing these tools can substantially increase your productivity
%	
%	\end{itemize}
%\end{frame}
%
%
%\begin{frame}{Utilize AI as Your Research Assistant}{Literature Review}
%	\begin{itemize} 
%		\item \textbf{ChatDOC}
%		
%		\vspace{0.2cm}
%			\begin{itemize} 
%				\item \url{https://chatdoc.com/}	
%								\vspace{0.2cm}
%		\item It can help you summarize the paper and answer any questions about the paper
%				\vspace{0.2cm}
%		\item You might still need to read the paper by yourself but this tool can help understand it better and save your time		
%	
%			\end{itemize}
%	\end{itemize}
%\end{frame}
%
%
%
%\begin{frame}{Utilize AI as Your Research Assistant}{Coding}
%	\begin{itemize} 
%		\item \textbf{ChatGPT}
%		
%		\vspace{0.2cm}
%		\begin{itemize} 
%			\item \url{https://chat.openai.com/auth/login}	
%			\vspace{0.2cm}
%			\item ChatGPT can help you draft the initial code 
%			\vspace{0.2cm}
%			\item It will explain every line in details so that you actually can learn how to coding from it
%			\vspace{0.2cm}
%			\item It can also help you debug or optimize your code
%			
%		\end{itemize}
%	\end{itemize}
%\end{frame}
%
%
%\begin{frame}{Utilize AI as Your Research Assistant}{Coding}
%	\begin{itemize} 
%		\item \textbf{ChatGPT}
%		
%		\vspace{0.2cm}
%		\begin{itemize} 
%			\item \url{https://res.cloudinary.com/dyd911kmh/image/upload/v1678453027/Marketing/Blog/ChatGPT_Cheat_Sheet.pdf}	
%			\vspace{0.2cm}
%			\item Whether ChatGPT can help you or not depends on how you ask it (give the informative prompts)  
%			\vspace{0.2cm}
%			\item The above cheat sheet will show some good examples 
%			
%			
%		\end{itemize}
%	\end{itemize}
%\end{frame}
%
%
%
%\begin{frame}{Utilize AI as Your Research Assistant}{Polish Your Writing}
%	\begin{itemize} 
%		\item \textbf{ChatGPT}
%		
%		\vspace{0.2cm}
%		\begin{itemize} 
%			\item You can use ChatGPT to polish your English writing
%					\vspace{0.2cm}
%			\item Rewrite or summarize your paragraph
%				\vspace{0.2cm}
%			\item It can also do a great job in translation		
%			
%		\end{itemize}
%	\end{itemize}
%\end{frame}
%
%
%\begin{frame}{Utilize AI as Your Research Assistant}{Polish Your Writing}
%	\begin{itemize} 
%		\item \textbf{Claude 3}
%		
%		\vspace{0.2cm}
%		\begin{itemize} 
%			\item \url{https://claude.ai/chats}
%			\vspace{0.2cm}
%			\item Claude 3 is also good at polishing and proofreading writing
%						\vspace{0.2cm}
%			\item It can also do coding very well			
%
%			
%		\end{itemize}
%	\end{itemize}
%\end{frame}



\end{document} 
